\section{Conclusiones}

Esta investigación analizó las disparidades socioeconómicas y la
convergencia regional de 24 departamentos del Perú entre 2007 y 2019, a
partir del Índice de Desarrollo Humano y del PBI real per cápita relativo.
La estrategia combinó estadística descriptiva, densidades kernel,
cartografía, autocorrelación espacial, regresiones globales, modelos de
error espacial, especificaciones Durbin y regresiones geográficamente
ponderadas.
 
El resultado más llamativo del estudio es también el que le da su título de
trabajo: la aproximación económica entre departamentos no vino acompañada de
una convergencia equivalente en desarrollo humano. El IDH promedio subió de
0.3488 a 0.5175 y su coeficiente de variación cayó de 0.1986 a 0.1624 —una
reducción de la dispersión relativa compatible con convergencia sigma—, pero
la regresión beta arrojó un coeficiente positivo y significativo de 0.1637.
Los departamentos que partían con mejores condiciones fueron, en promedio,
los que más avanzaron; no hay evidencia, entonces, de convergencia beta en
desarrollo humano, aun cuando la brecha relativa se haya estrechado.
 
El PBI real per cápita relativo cuenta una historia distinta. Su desviación
estándar bajó de 0.6193 a 0.4750 y el coeficiente inicial de la regresión
OLS fue de $-0.3282$, significativo por debajo del 1\%, con una velocidad
anual implícita cercana a 3.3\%. Los departamentos que arrancaban rezagados
en términos económicos mejoraron su posición relativa con mayor rapidez que
los mejor posicionados, aunque las diferencias de nivel entre ellos
siguieran siendo importantes al cierre del periodo.
 
Ni los niveles iniciales ni los cambios acumulados mostraron autocorrelación
espacial global significativa: los índices de Moran se mantuvieron cercanos
a cero, con valores de probabilidad por encima del 5\%, y esa conclusión no
varió al cambiar la matriz de distancia inversa por contigüidad queen o por
cuatro vecinos próximos en el caso del IDH. La relevancia del espacio, sin
embargo, no desaparece por eso: reaparece con fuerza distinta según el
indicador que se examine. En el desarrollo humano, los modelos espaciales no
alteran el signo positivo del coeficiente y la GWR arroja pendientes locales
casi idénticas entre departamentos. En el PBI per cápita, en cambio, los 24
coeficientes locales son negativos pero varían entre $-0.4857$ y $-0.2771$
—hay convergencia económica en todo el territorio, aunque con una intensidad
que cambia de un lugar a otro.
 
La elección de la matriz de vecindad tampoco es un detalle menor: el modelo
de error con contigüidad queen ofrece un mejor ajuste para el PBI per
cápita que el basado en banda de distancia inversa, el cual ni siquiera
mejora al OLS según el criterio de Akaike. Esto sugiere que cualquier
lectura de la dependencia espacial en este tipo de ejercicios debería
apoyarse en más de una definición de proximidad geográfica antes de sacar
conclusiones firmes.
 
En conjunto, la evidencia describe una trayectoria dual —acercamiento
económico junto con persistencia de diferencias en desarrollo humano— que
tiene implicancias directas para el diseño de política. Reducir las brechas
de ingreso relativo no garantiza, por sí solo, una distribución más
equilibrada de las capacidades sociales; las políticas regionales necesitan
integrar crecimiento productivo con educación, salud, infraestructura,
empleo y acceso a servicios. La velocidad desigual de la convergencia
económica, además, exige que las estrategias de desarrollo se adapten a la
estructura productiva, la conectividad y las condiciones iniciales de cada
departamento, sin perder de vista que los vínculos entre territorios
vecinos justifican una coordinación interdepartamental en infraestructura,
corredores económicos, mercados laborales y provisión de servicios.
 
Quedan varias líneas abiertas para investigaciones posteriores: paneles
anuales que capturen la dinámica dentro del periodo, unidades territoriales
más desagregadas, variables estructurales adicionales, y una exploración
específica de clubes de convergencia y dependencia espacial local más allá
de 2019. Avanzar en esa dirección ayudaría a precisar por qué el
acercamiento económico departamental no se traduce, todavía, en una
convergencia equivalente del desarrollo humano.
